\documentclass[12pt]{article}
\usepackage[utf8]{inputenc}
\usepackage[T1]{fontenc}
\usepackage[margin=1in]{geometry}
\usepackage{setspace}
\usepackage{xcolor}
\usepackage[colorlinks=true, urlcolor=blue!60!black, linkcolor=blue!60!black]{hyperref}
\setlength{\parindent}{2em}
\setlength{\parskip}{1.25em}
\renewcommand{\baselinestretch}{1.0}
\begin{document}
\begin{center}
{\LARGE \textit{Williamson v. Lee Optical of Oklahoma, Inc.}, 348 U.S. 483 (1955)\par}
\vspace{0.8em}
{\large Case Brief\par}
\end{center}
\vspace{1em}

\section{Memory Jogger}
Under the Fourteenth Amendment's Due Process and Equal Protection Clauses, economic regulations are upheld under hyper-deferential rational basis review whenever a court can conceive of any plausible legislative rationale---even if the law seems needless, wasteful, or addresses a problem only one step at a time.

\section{Detailed Case Facts}
In 1953, the Oklahoma Legislature enacted a comprehensive statutory scheme regulating visual care and the optical trade, codified at 59 Okla. Stat. Ann. \S\S~941--947 (Okla. Laws 1953, c. 13, \S\S~1--8) (\href{https://supreme.justia.com/cases/federal/us/348/483/}{Justia Opinion}). The optical industry involves three distinct roles: ophthalmologists (licensed physicians specializing in eye care), optometrists (licensed professionals who examine eyes for refractive error, recognize---but do not treat---eye diseases, and fill eyeglass prescriptions), and dispensing opticians (artisans trained to grind lenses, fill prescriptions, and fit frames to a customer's face) (\href{https://www.courtlistener.com/opinion/105285/williamson-v-lee-optical-of-oklahoma-inc/}{CourtListener}).

Four specific restrictions in the Oklahoma Act formed the core of the dispute (\href{https://supreme.justia.com/cases/federal/us/348/483/}{Justia Opinion}):
\begin{enumerate}
    \item \textbf{Prescription Requirement (\S~2):} Made it unlawful for anyone other than an Oklahoma-licensed ophthalmologist or optometrist to fit, adjust, or apply lenses, frames, or other optical appliances to a person's face, or to duplicate or replace lenses into frames, except upon the written prescriptive authority of a licensed Oklahoma ophthalmologist or optometrist. In practical effect, an optician could not place old lenses into new frames or duplicate a lost or broken lens using mechanical measuring devices without a written prescription.
    \item \textbf{Ready-to-Wear Exemption (\S~3):} Exempted sellers of ready-to-wear glasses equipped with convex-spherical lenses, plano sunglasses, and industrial safety goggles from the Act's regulatory framework when sold as merchandise at established places of business.
    \item \textbf{Advertising Ban (\S~3):} Prohibited soliciting the sale of spectacles, eyeglasses, lenses, frames, mountings, prisms, or other optical appliances, as well as eye examinations or visual services, via radio, television, window display, telephone directory display advertisements, or other advertising media.
    \item \textbf{Retail Space Restriction (\S~4):} Barred any retail merchandiser from renting space, subleasing departments, or otherwise permitting any person purporting to conduct eye examinations or provide visual care to occupy space inside a retail store.
\end{enumerate}

Lee Optical of Oklahoma, Inc., along with other dispensing opticians and affected businesses, sued Oklahoma Attorney General Mac Q. Williamson and other state officials seeking a declaratory judgment and injunction against enforcement of the statute under the Fourteenth Amendment's Due Process and Equal Protection Clauses (\href{https://www.oyez.org/cases/1940-1955/348us483}{Oyez}).

\section{Procedural History}
Because the plaintiffs sought to enjoin enforcement of a state statute on federal constitutional grounds under 28 U.S.C. \S\S~2201, 2202, and 2281, a three-judge panel of the U.S. District Court for the Western District of Oklahoma heard the case (\href{https://supreme.justia.com/cases/federal/us/348/483/}{Justia Opinion}). In \textit{Lee Optical of Oklahoma, Inc. v. Williamson}, 120 F. Supp. 128 (W.D. Okla. 1954), the District Court struck down parts of three sections of the Act:
\begin{itemize}
    \item It invalidated \S~2's prescription requirement for fitting old lenses into new frames and duplicating broken lenses under the Due Process Clause, finding that opticians could measure broken lenses with ordinary mechanical skills and that the means chosen were ``neither reasonably necessary nor reasonably related to the end sought to be achieved'' (120 F. Supp. at 137).
    \item It invalidated \S~3's exemption for ready-to-wear glasses under the Equal Protection Clause.
    \item It invalidated \S~3's ban on advertising eyeglass frames and mountings under the Due Process Clause, viewing frames as mere merchandise only casually related to public health (Judge Alfred P. Murrah dissented on this point below).
    \item It invalidated \S~4's ban on renting space in retail stores to eye-care providers under the Due Process Clause.
    \item However, the District Court upheld the portion of \S~3 barring advertisements for spectacles, eyeglasses, lenses, and prisms.
\end{itemize}

Both sides took direct appeals to the U.S. Supreme Court pursuant to 28 U.S.C. \S~1253: Oklahoma appealed the invalidations in No. 184, and Lee Optical cross-appealed the upheld advertising restriction in No. 185 (\href{https://supreme.justia.com/cases/federal/us/348/483/}{Justia Opinion}). The case was argued on March 2, 1955, and decided on March 28, 1955 (\href{https://www.oyez.org/cases/1940-1955/348us483}{Oyez}).

\section{Judicial Votes}
\textbf{Vote Breakdown:} 8--0 (Unanimous among participating Justices; 1 Justice not participating) (\href{https://www.oyez.org/cases/1940-1955/348us483}{Oyez}; \href{https://supreme.justia.com/cases/federal/us/348/483/}{Justia Opinion}).
\begin{itemize}
    \item \textbf{Majority (8):} Justice William O. Douglas (author), joined by Chief Justice Earl Warren and Justices Hugo L. Black, Stanley F. Reed, Felix Frankfurter, Harold H. Burton, Tom C. Clark, and Sherman Minton.
    \item \textbf{Concurrences (0):} None.
    \item \textbf{Dissents (0):} None.
    \item \textbf{Not Participating (1):} Justice John Marshall Harlan II took no part in the consideration or decision of the case (\href{https://supreme.justia.com/cases/federal/us/348/483/}{Justia Opinion}).
\end{itemize}

\section{Holding}
\textbf{Judgment:} Affirmed in part and reversed in part. The Supreme Court reversed the District Court's invalidation of \S\S~2, 3 (ready-to-wear exemption and frame-advertising ban), and 4, and affirmed the District Court's judgment upholding \S~3's ban on advertising spectacles, eyeglasses, lenses, and prisms (\href{https://supreme.justia.com/cases/federal/us/348/483/}{Justia Opinion}).

\textbf{Rule Announced:} Under the Fourteenth Amendment's Due Process Clause, a state economic or business regulation is constitutional so long as there is an evil at hand for correction and the legislature might rationally have thought the chosen measure was a reasonable way to correct it, regardless of whether the law is unwise, wasteful, or logically imperfect (\href{https://www.courtlistener.com/opinion/105285/williamson-v-lee-optical-of-oklahoma-inc/}{CourtListener}). Under the Equal Protection Clause, economic classifications violate the Constitution only if they amount to ``invidious discrimination''; a legislature may reform one phase of a field at a time while neglecting others (\href{https://supreme.justia.com/cases/federal/us/348/483/}{Justia Opinion}).

\textbf{Holding vs. Dicta:} The \textit{holding} is that none of the challenged provisions of Oklahoma's 1953 optical-care statute violated the Fourteenth Amendment's Due Process or Equal Protection Clauses because each bore a conceivable rational relationship to legitimate public health and professional-standards objectives. Justice Douglas's broader pronouncements---such as declaring that ``[t]he day is gone when this Court uses the Due Process Clause of the Fourteenth Amendment to strike down state laws, regulatory of business and industrial conditions,'' and quoting \textit{Munn v. Illinois}, 94 U.S. 113, 134 (1876), that ``[f]or protection against abuses by legislatures, the people must resort to the polls, not to the courts''---function as sweeping doctrinal dicta signaling the definitive burial of \textit{Lochner}-era economic substantive due process, though the Court did not formally overrule all pre-1937 cases or foreclose First Amendment commercial speech challenges to advertising bans that arose in later decades (\href{https://supreme.justia.com/cases/federal/us/348/483/}{Justia Opinion}).

\section{Analysis of Opinions}

\subsection{Majority Opinion (Justice Douglas)}
Justice Douglas systematically reversed each of the District Court's invalidations by supplying hypothetical rationales the Oklahoma Legislature ``might have concluded'' (\href{https://supreme.justia.com/cases/federal/us/348/483/}{Justia Opinion}):
\begin{enumerate}
    \item \textbf{Prescription Requirement (\S~2, Due Process):} Relying on \textit{Roschen v. Ward}, 279 U.S. 337 (1929), Douglas acknowledged that requiring a written prescription to replace frames or duplicate lenses ``may exact a needless, wasteful requirement in many cases.'' Nevertheless, balancing advantages and disadvantages belongs exclusively to the legislature. The legislature \textit{might} have concluded that prescriptions are necessary often enough to justify a blanket rule, or that eye examinations are so vital for detecting latent eye diseases that every frame change or lens duplication should prompt a visit to a specialist. Even though the statute did not actually require a new examination if an old prescription was on file, a law ``need not be in every respect logically consistent with its aims to be constitutional.'' Citing \textit{Nebbia v. New York}, 291 U.S. 502 (1934), and \textit{West Coast Hotel Co. v. Parrish}, 300 U.S. 379 (1937), Douglas declared the era of judicial invalidation of business regulations under substantive due process over (\href{https://supreme.justia.com/cases/federal/us/348/483/}{Justia Opinion}).
    \item \textbf{Ready-to-Wear Exemption (\S~3, Equal Protection):} Douglas emphasized that legislative classification admits of ``no doctrinaire definition.'' Reform may proceed ``one step at a time,'' addressing the phase of the problem that seems most acute to the legislature. The Equal Protection Clause forbids only ``invidious discrimination,'' and on the record, ready-to-wear glasses might not loom large in Oklahoma or might present distinct regulatory problems (\href{https://supreme.justia.com/cases/federal/us/348/483/}{Justia Opinion}).
    \item \textbf{Frame Advertising Ban (\S~3, Due Process):} Adopting Judge Murrah's dissent below, Douglas reasoned that an eyeglass frame is not used in isolation; it is used with lenses affecting human eyesight. The legislature might reasonably fear that frame advertisements would serve as bait to lure customers to buy lenses, and could therefore treat all who deal with the human eye as professionals subject to commercial advertising restrictions under \textit{Semler v. Oregon State Board of Dental Examiners}, 294 U.S. 608 (1935) (\href{https://supreme.justia.com/cases/federal/us/348/483/}{Justia Opinion}).
    \item \textbf{Retail Store Space Ban (\S~4, Due Process):} Analogizing to bans on corporate practice of dentistry in \textit{Semler}, Douglas held that barring eye-care practitioners from leasing space in retail stores was rationally related to freeing the profession from ``all taints of commercialism'' and preventing optometrists from becoming fronts for retail chains (\href{https://supreme.justia.com/cases/federal/us/348/483/}{Justia Opinion}).
    \item \textbf{Cross-Appeal on Lens Advertising (No. 185):} For the same reasons, the Court affirmed the District Court's ruling upholding \S~3's ban on advertising spectacles, eyeglasses, lenses, and prisms (\href{https://supreme.justia.com/cases/federal/us/348/483/}{Justia Opinion}).
\end{enumerate}

\subsection{Concurrences}
No Justice filed a separate concurring opinion (\href{https://supreme.justia.com/cases/federal/us/348/483/}{Justia Opinion}).

\subsection{Dissents}
No Justice filed a dissenting opinion (\href{https://supreme.justia.com/cases/federal/us/348/483/}{Justia Opinion}).

\section{Examples: Future Applications}
\begin{enumerate}
    \item \textbf{Same-Side (Upheld under \textit{Lee Optical}):} A state enacts a statute requiring all retail sellers of funeral caskets---even those who perform no embalming or burial services---to obtain a full funeral director's license requiring sixty hours of coursework and embalming training, while exempting online sellers located out of state. \textit{Analysis:} Under \textit{Lee Optical}, a court would uphold the law against a Due Process and Equal Protection challenge if it can hypothesize a conceivable consumer-protection rationale (e.g., ensuring grieving families interact with trained professionals), even if the law appears to be economic protectionism benefiting licensed funeral directors and leaves out-of-state sellers unregulated (one-step-at-a-time reform) (\href{https://ir.lawnet.fordham.edu/cgi/viewcontent.cgi?article=4853&context=flr}{Fordham Law Review}).
    \item \textbf{Same-Side (Upheld under \textit{Lee Optical}):} A state bans self-service gasoline stations and requires licensed attendants to pump all fuel, while allowing electric-vehicle drivers to plug in their own high-voltage chargers without an attendant. \textit{Analysis:} A legislature might rationally conclude that flammable liquid fuel poses distinct fire and spill hazards justifying trained attendants, and under \textit{Lee Optical}'s ``one step at a time'' Equal Protection doctrine, the state need not regulate every energy-refueling method simultaneously (\href{https://harvardlawreview.org/print/vol-138/when-rational-basis-review-bit/}{Harvard Law Review}).
    \item \textbf{Opposite-Side (Invalidated / Outside \textit{Lee Optical}):} A state passes a statute barring any person of a particular racial or national-origin group from obtaining an optician's license, claiming it streamlines administrative oversight. \textit{Analysis:} \textit{Lee Optical} applies only to ordinary economic classifications lacking ``invidious discrimination.'' Racial classifications trigger strict scrutiny under the Equal Protection Clause, where hypothetical rationales and deference do not apply.
    \item \textbf{Opposite-Side (Invalidated / Outside \textit{Lee Optical}):} A state statute prohibits licensed opticians from publishing truthful, non-misleading price lists for eyeglass frames on their websites, and the opticians challenge the law under the First Amendment rather than Fourteenth Amendment substantive due process. \textit{Analysis:} Although \textit{Lee Optical} upheld an advertising ban under Due Process review in 1955, modern First Amendment commercial speech doctrine subjects truthful commercial advertising bans to heightened intermediate scrutiny rather than \textit{Lee Optical} rational basis deference.
    \item \textbf{Fence-Sitter (Close / Contested Case):} A state requires independent African-style hair braiders---who use no chemicals, dyes, or scissors---to complete 1,500 hours of cosmetology school focused on chemical perms and hair cutting, and the state offers no evidence at trial that braiding poses any health risk. \textit{Analysis:} Under pure \textit{Lee Optical} deference, a court could imagine a tenuous sanitation rationale and uphold the requirement. Yet several federal courts applying ``rational basis with bite'' have questioned whether pure economic protectionism devoid of any factual link to public health can survive even rational basis review, creating a sharp doctrinal tension between \textit{Lee Optical}'s hypothetical-purpose test and anti-protectionism limits (\href{https://harvardlawreview.org/print/vol-138/when-rational-basis-review-bit/}{Harvard Law Review}; \href{https://ij.org/cje-post/the-rational-basis-test-the-story-continues/}{Institute for Justice}).
\end{enumerate}

\section{Critique}
\textbf{Originalist, Textualist, and Libertarian Critiques:} Libertarian and originalist scholars---including Randy Barnett and commentators at the Cato Institute, Pacific Legal Foundation, and Institute for Justice---widely criticize \textit{Lee Optical} as an abdication of the judicial duty to enforce constitutional limits on state power (\href{https://scholarship.law.georgetown.edu/cgi/viewcontent.cgi?referer=&httpsredir=1&article=2084&context=facpub}{Georgetown Law}; \href{https://www.cato.org/blog/hard-problem-abortion-rights}{Cato Institute}). They contend that the Fourteenth Amendment's Privileges or Immunities Clause and Due Process Clause were originally understood to protect the right to earn an honest living free from arbitrary ``class legislation'' and naked rent-seeking. By instructing judges to ignore the legislature's actual protectionist motives and invent hypothetical public-health rationales without requiring empirical evidence, \textit{Lee Optical} transformed rational basis review from a genuine check on arbitrary power into a rubber stamp for special-interest cartels (\href{https://ij.org/cje-post/the-rational-basis-test-the-story-continues/}{Institute for Justice}; \href{https://pacificlegal.org/discourse-irrational-basis/}{Pacific Legal Foundation}). A 2025 \textit{Harvard Law Review} Note similarly argues that \textit{Lee Optical} erred by abandoning the more balanced \textit{Carolene Products}-era rational basis standard, which had accorded legislatures deference while still policing naked economic protectionism through record-based means-end review (\href{https://harvardlawreview.org/print/vol-138/when-rational-basis-review-bit/}{Harvard Law Review}).

\textbf{Progressive and Doctrinal Perspectives:} Defenders of \textit{Lee Optical} from progressive and judicial-restraint traditions argue that the decision was essential to prevent courts from reviving \textit{Lochner v. New York}, 198 U.S. 45 (1905), and second-guessing democratic economic policy under the guise of constitutional law (\href{https://en.wikipedia.org/wiki/Lochner_v._New_York}{Wikipedia}). From this viewpoint, elected legislatures---not unelected federal judges---possess the institutional competence and democratic legitimacy to weigh competing economic interests, manage health and safety tradeoffs, and enact incremental reforms. At the same time, some progressive theorists note an irony: while \textit{Lee Optical}'s hyper-deference shields consumer and labor regulations from deregulation by conservative courts, it also leaves politically powerless low-income entrepreneurs vulnerable to entrenched industry lobbies when political processes are captured by economic power (\href{https://jackbalkin.yale.edu/footnote-part-iii}{Yale Law / Jack Balkin}).

\textbf{Underlying Competing Values:} The case squarely pits \textit{democratic self-governance and judicial restraint} (deferring to majoritarian policymaking on complex socioeconomic questions and avoiding judicial imposition of laissez-faire economics) against \textit{individual economic liberty and anti-capture constitutionalism} (preventing politically powerful incumbent guilds from using state police power to crush competitors and harm consumers).

\section{Key Quotations}

\begin{quote}
``The Oklahoma law may exact a needless, wasteful requirement in many cases. But it is for the legislature, not the courts, to balance the advantages and disadvantages of the new requirement.''
\end{quote}
\noindent --- Justice William O. Douglas, Majority Opinion, 348 U.S. at 487 (\href{https://supreme.justia.com/cases/federal/us/348/483/}{Justia}; \href{https://www.courtlistener.com/opinion/105285/williamson-v-lee-optical-of-oklahoma-inc/}{CourtListener}).

\begin{quote}
``But the law need not be in every respect logically consistent with its aims to be constitutional. It is enough that there is an evil at hand for correction, and that it might be thought that the particular legislative measure was a rational way to correct it.''
\end{quote}
\noindent --- Justice William O. Douglas, Majority Opinion, 348 U.S. at 487--488 (\href{https://supreme.justia.com/cases/federal/us/348/483/}{Justia}; \href{https://www.courtlistener.com/opinion/105285/williamson-v-lee-optical-of-oklahoma-inc/}{CourtListener}).

\begin{quote}
``The day is gone when this Court uses the Due Process Clause of the Fourteenth Amendment to strike down state laws, regulatory of business and industrial conditions because they may be unwise, improvident, or out of harmony with a particular school of thought.''
\end{quote}
\noindent --- Justice William O. Douglas, Majority Opinion, 348 U.S. at 488 (\href{https://supreme.justia.com/cases/federal/us/348/483/}{Justia}; \href{https://www.courtlistener.com/opinion/105285/williamson-v-lee-optical-of-oklahoma-inc/}{CourtListener}).

\begin{quote}
``For protection against abuses by legislatures, the people must resort to the polls, not to the courts.''
\end{quote}
\noindent --- Justice William O. Douglas, Majority Opinion, 348 U.S. at 488 (quoting \textit{Munn v. Illinois}, 94 U.S. 113, 134 (1876)) (\href{https://supreme.justia.com/cases/federal/us/348/483/}{Justia}; \href{https://www.courtlistener.com/opinion/105285/williamson-v-lee-optical-of-oklahoma-inc/}{CourtListener}).

\begin{quote}
``The problem of legislative classification is a perennial one, admitting of no doctrinaire definition. Evils in the same field may be of different dimensions and proportions, requiring different remedies. Or so the legislature may think.''
\end{quote}
\noindent --- Justice William O. Douglas, Majority Opinion, 348 U.S. at 489 (\href{https://supreme.justia.com/cases/federal/us/348/483/}{Justia}; \href{https://www.courtlistener.com/opinion/105285/williamson-v-lee-optical-of-oklahoma-inc/}{CourtListener}).

\end{document}
